\chapter*{Course Map: From a Program to a Complete Operating System}
\addcontentsline{toc}{chapter}{Course Map: From a Program to a Complete Operating System}
\markboth{Course Map}{}

An operating system sits between application programs and hardware. The course becomes easier when every topic is attached to one of four questions: \emph{what runs, who gets the CPU, where does data live in memory, and how does persistent data reach storage?}

\begin{mlfigure}
\begin{tikzpicture}[
  node distance=0.55cm,
  stage/.style={draw=MLBlue,fill=MLPaleBlue,rounded corners=2mm,text width=12.2cm,minimum height=1.25cm,align=left,inner xsep=4mm,inner ysep=2.5mm,font=\small\color{MLNavy}},
  arr/.style={-{Stealth[length=2.3mm]},thick,draw=MLTeal}
]
\node[stage] (s1) {\textbf{1. Interface to the OS} --- Linux commands, files, directories, processes, \texttt{fork()}, and basic system interfaces.};
\node[stage,below=of s1] (s2) {\textbf{2. CPU and concurrency} --- scheduling, process interaction, semaphores, shared memory, POSIX threads, mutexes, and deadlock.};
\node[stage,below=of s2] (s3) {\textbf{3. Memory} --- contiguous allocation, fragmentation, virtual memory, page faults, and replacement policies.};
\node[stage,below=of s3] (s4) {\textbf{4. Persistent storage} --- directory trees, file records, file-system metadata, and indexed allocation.};
\draw[arr] (s1) -- (s2); \draw[arr] (s2) -- (s3); \draw[arr] (s3) -- (s4);
\end{tikzpicture}
\end{mlfigure}

\begin{firstread}
Use the labs as checkpoints rather than isolated programs. Each one exposes a real OS abstraction: a process identifier, an open directory stream, a scheduling queue, a semaphore count, a safe-state test, a memory hole, a page frame, or an indexed block.
\end{firstread}

\begin{mltable}
\begin{tabularx}{\linewidth}{@{}>{\bfseries\leavevmode\color{MLNavy}}p{0.16\linewidth}p{0.30\linewidth}X@{}}
\toprule
\rowcolor{MLPaleBlue!92}
Lab & Practical focus & Core OS idea \\
\midrule
1 & Directories and \texttt{fork()} & OS services, processes, PIDs \\
2 & File text processing and directory reading & File I/O and directory APIs \\
3 & Bash scripting & Command interpreter and automation \\
4 & FCFS, SJF, Priority, Round Robin & CPU scheduling \\
5 & Producer-consumer and shared memory & Synchronization and IPC \\
6 & Banker's algorithm, deadlock detection, threads & Resource safety and concurrency \\
7 & First, best, and worst fit & Contiguous memory allocation \\
8 & Mutexes, FIFO, LRU, LFU & Mutual exclusion and page replacement \\
9 & Directory visualization & Hierarchical file-system namespace \\
10 & Records and indexed allocation & File organization and block allocation \\
\bottomrule
\end{tabularx}
\end{mltable}

\begin{keyidea}
Do not memorize the algorithms as unrelated loops. For every algorithm, ask: \textbf{what state does the OS maintain, what event changes that state, and what policy decides the next action?}
\end{keyidea}
